Este informe corrió ocho campañas de barrido dedicadas en \texttt{paccaA100},
separadas por completo del árbol de resultados de las campañas reales de
construcción del catálogo (\texttt{hyperion-results/sweeps/} vs.\
\texttt{hyperion-results/campaigns/}), para auditar con evidencia directa
los parámetros de diseño experimental de la Fase 1. La
Tabla~\ref{tab:run-accounting} desglosa la contabilidad exacta de corridas,
separando \textbf{corridas válidas retenidas para el análisis} de
\textbf{invocaciones inválidas ya reemplazadas} -- estas últimas no se
descartan del registro histórico: son precisamente la evidencia de los dos
defectos de infraestructura documentados en este informe
(\texttt{dgemm\_n2048}/\texttt{libopenblas.so.0}, y el video truncado de
\texttt{rodinia\_heartwall}).

\begin{table}[H]
\centering
\caption{Contabilidad de invocaciones del harness, por barrido.}
\label{tab:run-accounting}
\small
\begin{tabular}{@{}lrrl@{}}
\toprule
\textbf{Barrido} & \textbf{Válidas} & \textbf{Inválidas reemplazadas} & \textbf{Cálculo (válidas)} \\
\midrule
\texttt{interval\_ns} (Sección~\ref{sec:interval-ns}) & 105 & 15 (\texttt{dgemm\_n2048}) & 7 kernels $\times$ 5 cadencias $\times$ 3 reps \\
\texttt{gpu\_interval\_ns} (Sección~\ref{sec:gpu-interval-ns}) & 120 & 15 (\texttt{rodinia\_heartwall}) & 8 kernels $\times$ 5 cadencias $\times$ 3 reps \\
\texttt{repetitions} (Sección~\ref{sec:repetitions}) & 150 & 20 (10 \texttt{dgemm\_n2048} + 10 \texttt{rodinia\_heartwall}) & 15 kernels $\times$ 10 reps \\
\texttt{smt\_policy} (Sección~\ref{sec:smt-policy}) & 10 & 0 & 2 políticas $\times$ 5 reps \\
Confirmación aleatorizada (Sección~\ref{sec:randomized-confirmation}) & 24 & 0 & 18 CPU + 6 GPU \\
\midrule
\textbf{Total} & \textbf{409} & \textbf{50} & \\
\bottomrule
\end{tabular}
\end{table}

En total, el harness se invocó \textbf{459 veces} (409 válidas + 50
inválidas reemplazadas) a lo largo de estos ocho barridos. Todas las
cifras de este informe (medias, CV\%, tablas) usan únicamente las 409
corridas válidas -- las 50 inválidas se muestran aquí por transparencia de
auditoría, no porque contribuyan al análisis. Aparte de estas, se
perfilaron 24 combinaciones independientes con
\texttt{ncu} (21 del barrido F original + 3 de la extensión de
\texttt{rodinia\_lud}) y se corrieron 50 calibraciones repetidas -- ninguna
de las dos usa \texttt{runner.run\_single()}, así que no se suman al total
de corridas del harness. Los reintentos por fallos de infraestructura
(\texttt{dgemm\_n2048} por \texttt{libopenblas.so.0} ausente,
\texttt{rodinia\_heartwall} por el video truncado) \textbf{reemplazan}
corridas ya contadas en su barrido original -- no se cuentan aparte, para
no inflar el total con corridas que en su momento fueron inválidas y
después se repitieron. La Sección~\ref{sec:master-table} recoge el estado
real de los 24 parámetros identificados originalmente; este balance resume
las conclusiones cualitativas, con el cuidado de no reclamar más de lo que
cada medición sostiene.

\textbf{Valor razonable y conservador, verificado en esta plataforma} --
la formulación correcta para la mayoría de estos parámetros, distinta de
``valor óptimo demostrado'': \texttt{interval\_ns} (1\,ms es el primer
punto donde el riesgo de jitter severo cae de forma abrupta y el IPC ya es
estable, sin que esto excluya que otros valores cercanos también serían
válidos, Sección~\ref{sec:interval-ns}); \texttt{gpu\_interval\_ns} (5\,ms
es viable y conservador, pero la resolución real del sensor NVML
($\sim$100--120\,ms) hace que el margen efectivo sobre el objetivo de
ventanas sea mucho más estrecho de lo que sugiere contar lecturas crudas,
y 10\,ms también sería viable, Sección~\ref{sec:gpu-interval-ns});
\texttt{repetitions\_per\_combination} (3 es un mínimo operativo razonable
para Fase 1, no una garantía estadística -- el reanálisis combinatorio
completo sobre las 120 ternas posibles de cada kernel muestra que el
resultado depende de qué terna se ejecute, y que la limitación es más
seria para la futura evaluación energética de Fase 4,
Sección~\ref{sec:repetitions}); y \texttt{target\_windows\_per\_repetition}
(un guardarraíl de duración mínima nunca activo en 150 corridas, no una
prueba de suficiencia estadística de la muestra resultante).

\textbf{Confirmado con medición directa, sin matices adicionales}: el
número de lanzamientos de \texttt{ncu} por kernel (con un criterio de
convergencia declarado de antemano, y con la extensión dirigida que cierra
el caso más exigente, \texttt{rodinia\_lud}, en 200 lanzamientos, por
encima de su punto de convergencia real de 150,
Sección~\ref{sec:ncu-convergence}).

\textbf{Efecto observado con medición directa, causa y ventaja energética
no aisladas}: \texttt{smt\_policy} -- en \texttt{npb\_cg}, la
configuración de 12 hilos con SMT completo reduce el IPC agregado en
40{,}4\,\% frente a 6 hilos en 6 núcleos físicos, un efecto sistemático y
no ruido de medición (Sección~\ref{sec:smt-policy}). Pero el diseño cambió
\texttt{OMP\_NUM\_THREADS} (6$\to$12) a la vez que la política de SMT, así
que MPKI y tasa de fallos de LLC constantes descartan más tráfico de caché
por unidad de acceso, pero no descartan mayor tráfico total por unidad de
tiempo, ni aíslan contención de \textit{pipeline} de otras causas
(sincronización de OpenMP, ancho de banda de memoria). Tampoco se midió
energía, y el tiempo de pared fue menor bajo SMT completo -- así que este
resultado no permite afirmar que \texttt{one\_thread\_per\_physical\_core}
sea la política energéticamente superior. Se mantiene por consistencia de
telemetría en Fase 1, no como óptimo de EDP confirmado.

\textbf{Confirmado en dirección, no en magnitud extrema}: la confirmación
aleatorizada de los dos casos frontera (Sección~\ref{sec:randomized-confirmation})
tiene fuerza distinta en cada dominio. En GPU, \texttt{rodinia\_backprop}
reprodujo exactamente 17 y 9 muestras bajo orden aleatorizado -- una
confirmación fuerte, cifra por cifra. En CPU, la confirmación solo
reprodujo una diferencia moderada en la media de CV\% entre 0{,}5\,ms y
1\,ms para 3 kernels -- no reapareció la cola severa (el máximo de
8{,}448\,\% que es el argumento central contra 0{,}5\,ms en la
Sección~\ref{sec:interval-ns}), porque esa cola depende de un evento
ocasional de interferencia del planificador que 3 repeticiones por punto
no garantizan capturar. La confirmación aleatorizada de CPU respalda la
\emph{dirección} del hallazgo (0{,}5\,ms es más ruidoso), no una
reconfirmación del riesgo de cola extrema en sí.

\textbf{Guardarraíl conservador confirmado, no tolerancia derivada
estadísticamente}: \texttt{D03\_TOLERANCE\_FRACTION} (40\,\%, CPU) -- el
ruido real medido está casi tres órdenes de magnitud por debajo, pero eso
confirma que el guardarraíl nunca se ha acercado a su límite, no que 40\,\%
sea el valor óptimo frente a 10\,\% o 20\,\%. La tolerancia cruzada de GPU
(5\,\%) queda \textbf{parcialmente justificada}: tiene el mismo margen
frente al ruido, pero su capacidad real de detectar una aplicación fallida
de frecuencia sigue sin poder verificarse hasta el permiso P4
(Sección~\ref{sec:calibration-distribution}).

\textbf{Justificados por decisión metodológica explícita, no por
medición} (Sección~\ref{sec:literatura}): los niveles de frecuencia F0-F4
(sin permiso de escritura disponible todavía) y el umbral de carga externa
E08 (razonable por convención de sistemas). El rango de temperatura E02
queda \textbf{parcial y corregido}: la hoja de especificaciones exacta del
SKU del nodo (Xeon Gold 5315Y) documenta un $T_{\text{case}}$ máximo de
81\,$^{\circ}$C, no 90\,$^{\circ}$C, así que 90\,$^{\circ}$C debe
presentarse como margen de seguridad conservador frente al rango típico de
$T_j$máx de la familia, no como un umbral verificado contra el datasheet
exacto de este procesador -- una corrección respecto a una afirmación más
fuerte en una versión anterior de este análisis.

\textbf{Hallazgo de integración, más serio que la falta de barrido de
E08/E02}: revisando el código directamente (no solo el diseño conceptual),
se encontró que \textbf{ni E08 ni E02 se ejecutan hoy en ninguna campaña
real del proyecto}, y que el preflight de GPU (G01--G03) tampoco corre
porque el manifiesto de referencia GPU declara \texttt{gpu.enabled: false}.
\texttt{run\_campaign\_preflight} (invocado una vez, al inicio de la
campaña, desde \texttt{cli.py}) no incluye E02 ni E08; ambos viven en
\texttt{run\_reduced\_preflight}, pensado para correr antes de cada
combinación, pero que \texttt{campaign.py} nunca llama (sí verifica E06,
procesos ajenos, antes de cada combinación -- pero no E02 ni E08). Y
aunque se llamara, \texttt{Manifest} no declara
\texttt{package\_temperature\_c}, así que E02 aprobaría siempre por sensor
ausente; y aunque \texttt{gpu.enabled} fuera \texttt{true}, \texttt{cli.py}
no le pasa ningún inspector NVML a G01--G03, así que tampoco podrían
verificar procesos CUDA ajenos, modo persistente ni configuración MIG.
Esto no es una limitación del alcance de este informe -- es un hallazgo
sobre el instrumento de producción mismo, documentado en detalle en la
Sección~\ref{sec:metodologia} y en la Sección~\ref{sec:literatura}, con la
misma seriedad que el defecto de datos de \texttt{rodinia\_heartwall}
descrito más abajo. Por alcance de un proyecto de pregrado, este informe
documenta el hallazgo y lo deja como prerrequisito explícito antes de la
campaña DVFS definitiva -- integrar estos tres chequeos en
\texttt{campaign.py} es trabajo de implementación fuera del alcance de
esta auditoría de parámetros, no algo que este informe deba resolver.

\textbf{Fijado por diseño para evitar sesgo de selección posterior}: el
peso $w$ del EDP se pre-registra en $w=1$ (EDP clásico) como métrica
primaria de Fase 4, con $w=2$ ($ED^2P$) como análisis de sensibilidad
secundario -- decidir esto \emph{antes} de tener resultados de Fase 4 es
lo que evita, aunque no fuera la intención original, escoger
implícitamente la métrica que favorezca al agente de control
(Sección~\ref{sec:literatura}).

\textbf{Limitación metodológica declarada, no oculta}: los barridos
principales corrieron sus combinaciones en orden fijo, no aleatorizado,
sin preflight repetido de carga externa/temperatura por combinación -- un
riesgo real en un nodo compartido. Una confirmación dirigida y aleatorizada
de los dos casos frontera más citados reprodujo el mismo resultado
(Sección~\ref{sec:randomized-confirmation}), pero eso acota el riesgo solo
para esos dos casos, no para la totalidad de los barridos.

\textbf{Hallazgo no buscado, con impacto real fuera del alcance original
de este informe}: se encontró y corrigió un defecto de datos de producción
-- \texttt{rodinia\_heartwall} corría, desde algún momento posterior a su
caracterización original (ARC-86), contra un video sintético truncado a 25
cuadros en vez de los 20000 esperados, sin que el harness lo detectara
(el binario de Rodinia termina con código de salida 0 tras su propio error
de validación de argumentos). Esto invalidaba silenciosamente tres
repeticiones ya marcadas \texttt{accepted} en la campaña de producción y
los resultados de \texttt{rodinia\_heartwall} en dos de los barridos de
este mismo informe. Se corrigió la causa raíz, se recuperaron los datos de
producción, y se re-corrieron los barridos afectados. Este episodio es, en
sí mismo, la mejor evidencia del valor de este ejercicio de auditoría: no
se estaba buscando este problema, y no se habría encontrado sin medir de
verdad en vez de asumir que los datos ya existentes eran correctos. La
misma disciplina de verificación llevó, en una revisión posterior de este
informe, a corregir una segunda afirmación demasiado fuerte sobre
\texttt{rodinia\_heartwall} y sobre el rango de temperatura E02, y a
descubrir que la resolución real del sensor NVML es mucho más gruesa de lo
que el conteo ingenuo de lecturas sugería.

\textbf{Alcance y límites del informe, dichos sin rodeos}: este trabajo
valida que la instrumentación y varios parámetros de adquisición funcionan
de forma reproducible en \texttt{paccaA100}, y justifica decisiones
locales de ingeniería para construir el dataset de Fase 1. No demuestra
que estos parámetros sean universales para otras plataformas. No valida
todavía el efecto de DVFS, el ahorro energético ni el EDP final, porque
faltan las campañas reales multifrecuencia (bloqueadas por los permisos
P1/P4). Tampoco convierte automáticamente las ventanas recolectadas en un
dataset estadísticamente suficiente para entrenar un modelo -- eso deberá
verificarse en Fase 2 con partición por corrida/carga, balance de clases y
evaluación fuera de muestra, un análisis que este informe no realiza.

De los 24 parámetros originalmente identificados, la clasificación
canónica es la de la Tabla~\ref{tab:master}
(Sección~\ref{sec:master-table}) -- una versión anterior de este balance
daba un conteo distinto (12/3/4) que no coincidía con esa tabla; se
corrige aquí para que ambos sean la misma fuente. Ninguno de los 24 queda
sin veredicto explícito: \textbf{10} quedan medidos con barrido dedicado
en este informe y sus extensiones sin matices adicionales, \textbf{2}
provienen de medición directa ya reportada en el documento anterior,
\textbf{5} quedan parcialmente justificados con una condición de cierre
concreta (tolerancia cruzada GPU, \texttt{MARGIN}, \texttt{plateau\_ratio},
rango de temperatura E02, y \texttt{delegated\_cpus}), \textbf{2} quedan
justificados por decisión metodológica o convención explícita (F0-F4, E08),
\textbf{1} queda fijado por diseño para evitar sesgo de selección
posterior (peso del EDP), y \textbf{4} permanecen explícitamente sin
justificar, sin haber sido tocados por ningún barrido hasta ahora. La
diferencia frente a la primera versión de este informe no es cuántos
parámetros tienen veredicto -- ya eran todos -- sino qué tan fuerte es la
palabra que describe cada veredicto: ``razonable y verificado en esta
plataforma'' en vez de ``óptimo'' o ``confirmado'' donde la evidencia no
alcanza para lo segundo.
